%=====================================================================
%  02 --- SYSTEM DESCRIPTION
%  Team 1 --- Uber | Delivery 1: Requirements & Workload Model
%  Owner: Jhonatan Barragán
%
%  Covers: Business Problem, Actors, Main Use Cases, System Boundaries
%  (+ system boundary diagram).
%
%  REQUIRES in main.tex preamble (add once, does not affect other
%  sections):
%      \usepackage{tikz}
%      \usetikzlibrary{arrows.meta}
%  Citations \cite{...} resolve against the entries added to
%  06-references.tex.
%=====================================================================
\section{System Description}
\label{sec:system-description}

This section sets out what the platform has to do and the conditions it has
to work under, before any architectural or technological choice is made. As
the delivery requires, it stays technology-agnostic: no cloud provider,
database, or implementation choice is assumed. We begin with the business
problem that motivates the system and why it becomes hard at scale, then move
to the actors that interact with it, the main use cases those actors drive,
and finally the boundary that separates the system from everything around it.

%---------------------------------------------------------------------
\subsection{Business Problem}
\label{subsec:business-problem}

\subsubsection*{The pain point}

City transportation has long carried a coordination failure between two
groups that never see each other in time: people who need a ride at a given
moment, and drivers with idle capacity at that same moment. Under the older
model of street hailing and phone-dispatched taxis, neither side has reliable
real-time information about the other. A rider cannot tell whether a car is
nearby, how long it will take, or what the trip will cost until it is over; a
driver has no clear picture of where unmet demand sits and ends up spending a
good part of the day driving empty. The visible symptoms are long and
unpredictable waits, poor vehicle use, opaque fares, cash-only friction, and
little accountability on either safety or service.

Uber's answer is to run a two-sided marketplace that connects ride demand
with independent driver supply in real time, without owning the cars. The
platform closes the information gap between the two sides: it finds nearby
drivers, shows an up-front price and an estimated arrival time, coordinates
the trip from request to drop-off, and settles payment on its own. An
informal, badly coordinated activity becomes a measurable on-demand service.

\subsubsection*{Who benefits}

The value is different for each side of the market and for the environment
around it. \textbf{Riders} get on-demand availability, predictable ETAs,
transparent and cashless pricing (or cash, still common in much of Latin
America), trip tracking, and safety mechanisms such as identity checks and a
shareable trip status. \textbf{Drivers} get flexible earning opportunities
and a demand signal that tells them where and when riders are requesting,
which raises how much of their time and vehicle actually earns.
\textbf{Cities and the wider economy} gain a mobility option that complements
public transit and can reduce the need to own a car for occasional trips.

The scale of this coordination is large. Uber reported 11.273 billion trips
and \$162.8 billion in gross bookings for fiscal year 2024, and in the fourth
quarter of 2024 it served 171 million monthly active platform consumers, an
average of roughly 33 million trips a day~\cite{uber2024fy}. The platform
runs in around 70 countries and more than 10{,}000 cities~\cite{uberH3}. Each
of those trips is a separate, geographically local coordination event that
has to resolve in seconds.

\subsubsection*{Why it is hard at scale}

What makes Uber a genuine distributed-systems problem is not the raw volume
on its own, but the set of constraints that have to hold at the same time for
every one of those events:

\begin{enumerate}[leftmargin=1.6em,itemsep=0.2em]
  \item \textbf{Perishable, real-time inventory.} The unit of supply is a
        driver-minute, and it cannot be stored: a minute a driver spends idle
        is gone. The market therefore has to clear continuously and in real
        time, not in batches, so that supply and demand meet while both still
        exist.

  \item \textbf{Local decisions on a global system.} Matching is inherently
        local, since it concerns drivers within a few kilometers of a rider,
        yet these local decisions happen at once across thousands of cities.
        The system has to reason at fine geographic granularity and at very
        large aggregate volume. Uber's H3 hexagonal index was built to bucket
        marketplace events by area so that supply, demand, and matching can be
        analyzed per zone in every city it serves~\cite{uberH3}.

  \item \textbf{Latency as a product requirement.} A rider expects a nearby
        driver and an ETA within seconds of requesting, and driver locations
        stream in constantly and must stay fresh. How fast matching responds
        is part of the product, not an implementation detail.

  \item \textbf{Dynamic market clearing.} When demand in an area outruns
        supply, price has to rise to rebalance the two sides, attracting more
        drivers and tempering demand, and it has to do so per area and in
        near real time. Uber derives this surge by measuring supply and
        demand within hexagonal cells in each city~\cite{uberH3}.

  \item \textbf{Bursty, uneven load.} Demand follows time of day, weather,
        and events, so peaks (rush hours, concerts, New Year's Eve) land far
        above the daily average and in particular places. The core matching
        experience has to hold up through those spikes.

  \item \textbf{Different correctness needs in one request.} Parts of the
        platform sit at different points of the consistency--availability
        trade-off. Payments, driver payouts, and trip state need strong
        correctness and a clear audit trail; money cannot be lost or charged
        twice. Location updates and driver discovery, on the other hand,
        favor low latency and high availability and can live with weaker,
        eventual consistency. A single request often touches both at once.

  \item \textbf{Trust and safety as requirements, not features.} The platform
        puts strangers together and moves real money at scale, so identity
        verification, ratings, fraud prevention, and trip safety are
        load-bearing requirements that shape the whole design.

  \item \textbf{Geographic and regulatory spread.} Users sit across many
        countries and rule sets, which brings both latency constraints
        (serving users near where they are) and functional variation (pricing
        rules, driver requirements, and payment methods differ by market).
\end{enumerate}

Put together, these constraints frame the core engineering problem: keep
clearing a perishable, two-sided, local marketplace, in real time and at
global scale, while holding strong correctness for money and safety and high
availability for matching. The requirements in the rest of this document, and
the architecture decisions in the deliveries that follow, all trace back to
this.

%---------------------------------------------------------------------
\subsection{Actors}
\label{subsec:actors}

An actor is any entity, human or external system, that sits outside the
system's boundary and interacts with it to reach a goal. Two things follow
from that definition. First, parties that are affected by the platform but
never interact with it directly (cities, regulators, and the wider economy
from Section~\ref{subsec:business-problem}) are indirect stakeholders, not
actors, and do not appear here. Second, components that live inside the
system (the matching, pricing, or fraud logic, for example) are not actors
either; only entities the system does not own qualify.

We classify actors on two axes. By nature they are either human or system
(external) actors. By role in an interaction, a \emph{primary} actor starts a
use case for its own goal, while a \emph{secondary} actor is called on by the
system to help finish one. For each actor we give its role, its typical
interaction frequency, and what it expects of the system. Since this is a
requirements-level view of the global platform, frequency is described in
qualitative terms; the numbers behind it are left to the non-functional
requirements and the workload model.

\subsubsection*{Human actors}

\begin{itemize}[leftmargin=1.4em,itemsep=0.2em]
  \item \textbf{Rider (Passenger) --- primary.} The demand side. The rider
        requests a trip, is matched with a driver, travels, pays, and rates
        the trip. The interaction is occasional but active and
        latency-sensitive over the request--trip--payment cycle. A rider
        expects a nearby driver within seconds, an accurate ETA and up-front
        price, a safe and trackable ride, easy payment (card or cash), and a
        real answer when something goes wrong.

  \item \textbf{Driver --- primary.} The supply side. The driver goes online
        to offer availability, receives and accepts offers, drives to and
        carries the rider, and gets paid. While online the interaction is
        continuous: location streams constantly and offers arrive in real
        time. A driver expects a steady flow of nearby requests, transparent
        and timely earnings, workable navigation, a sense of where demand is,
        and support when needed.

  \item \textbf{Customer Support / Safety Agent --- secondary.} A person who
        handles incidents, disputes, account problems, and safety
        escalations, including live help during a trip. The interaction is
        event-driven, triggered only when a rider or driver raises something,
        but it has to be available around the clock. The agent needs access
        to trip history and live trip state, plus the tools to act (refunds,
        account changes, escalation).

  \item \textbf{Operations / Market Administrator --- secondary.} An internal
        operator responsible for a city or region: reviewing driver
        onboarding, setting pricing and service-area policy, and watching
        marketplace health. This is administrative and periodic work, not
        per-trip. The administrator needs dashboards and controls, the
        ability to set policy per market, and visibility into supply, demand,
        and service quality.
\end{itemize}

\subsubsection*{System (external) actors}

\begin{itemize}[leftmargin=1.4em,itemsep=0.2em]
  \item \textbf{Payment Service Provider --- secondary.} Authorizes and
        captures rider charges and pays drivers out through the supported
        payment methods. It is called once per trip (authorization at the
        start, capture at completion) plus on the payout schedule. From it
        the system needs secure, reliable, idempotent transactions and a
        firm settlement result.

  \item \textbf{Mapping \& Geospatial Data Provider --- secondary.} Supplies
        the base map data and geocoding used to place riders and drivers,
        estimate ETAs, and drive navigation. Much of the routing and
        geospatial logic is built in-house and sits inside the boundary; the
        actor here is only the external base-map and geodata source. It is
        used continuously and at high volume, and the system needs accurate,
        low-latency responses from it.

  \item \textbf{Identity \& Background Verification Provider --- secondary.}
        Verifies driver identity and runs background screening at onboarding,
        with periodic re-screening and real-time identity checks
        thereafter~\cite{uberScreening}. It is used mostly at onboarding and
        on the re-check schedule, and the system needs authoritative,
        auditable results from it.

  \item \textbf{Notification / Communications Provider --- secondary.}
        Delivers one-time passwords, push notifications, SMS, and email (trip
        updates, receipts). It is used often and is event-driven, with
        several messages per trip, and what the system needs from it is
        timely, dependable delivery.
\end{itemize}

\subsubsection*{Actor catalog summary}

Table~\ref{tab:actors} gathers the catalog for quick reference and for
traceability in later deliveries.

\begin{table}[H]
  \centering
  \small
  \setlength{\tabcolsep}{4pt}
  \renewcommand{\arraystretch}{1.3}
  \begin{tabular}{@{}p{3.0cm} p{1.4cm} p{2.5cm} p{5.0cm}@{}}
    \toprule
    \textbf{Actor} & \textbf{Nature} & \textbf{Frequency} &
    \textbf{Key expectations} \\
    \midrule
    Rider (Passenger) & Human & Occasional, active &
      Fast matching, accurate ETA and price, safe trackable ride,
      easy payment, support \\
    Driver & Human & Continuous while online &
      Steady nearby requests, transparent fair earnings, navigation,
      demand signal, safety \\
    Support / Safety Agent & Human & Event-driven (24/7) &
      Access to trip data and tools to resolve issues and escalate
      safety events \\
    Operations / Market Admin & Human & Periodic, administrative &
      Dashboards, per-market policy controls, marketplace visibility \\
    Payment Service Provider & System & Per trip &
      Secure, reliable, idempotent transactions and settlement
      confirmation \\
    Mapping \& Geospatial Data & System & Continuous, high volume &
      Accurate, low-latency map, geocoding, and ETA data \\
    Identity \& Background Verif. & System & Onboarding + periodic &
      Authoritative, auditable identity and screening results \\
    Notification / Comms & System & Event-driven, frequent &
      Timely, reliable delivery of messages and one-time passwords \\
    \bottomrule
  \end{tabular}
  \caption{Actor catalog for the Uber platform (global scope).}
  \label{tab:actors}
\end{table}

%---------------------------------------------------------------------
\subsection{Main Use Cases}
\label{subsec:use-cases}

This subsection specifies the platform's primary use cases, the ones that
carry its core value. Supporting, exceptional, and purely administrative
interactions are left out to keep the model focused. Each use case follows
the same template: the actors involved (from the catalog in
Section~\ref{subsec:actors}), preconditions, the main (success) flow, the
relevant alternative flows, and postconditions. Each carries an identifier
(UC-\emph{nn}) so that functional requirements in later sections can trace
back to it.

\subsubsection*{UC-01 --- Request and match a ride}

\begin{description}[leftmargin=2.4cm,style=nextline,font=\normalfont\bfseries]
  \item[Primary actor] Rider (Passenger).
  \item[Secondary actors] Driver (supply side); Mapping \& Geospatial Data
        Provider.
  \item[Preconditions] The rider is authenticated with a valid account and a
        registered payment method; the rider's location is available; the
        request channel is operational.
  \item[Main flow]
    \begin{enumerate}[leftmargin=1.6em,itemsep=0.15em]
      \item The rider specifies pickup and destination.
      \item The system estimates the route, ETA, and up-front price for the
            available service tiers.
      \item The rider reviews the offer and confirms the request.
      \item The system finds nearby available drivers and dispatches the
            request according to its matching policy.
      \item A driver accepts the request.
      \item The system confirms the match to the rider, showing driver and
            vehicle details and the arrival ETA.
    \end{enumerate}
  \item[Alternative flows]
    \begin{itemize}[leftmargin=1.6em,itemsep=0.15em]
      \item \textbf{A1 --- No supply available.} No eligible driver is found
            nearby; the system tells the rider and may offer to keep
            searching or to notify when supply appears.
      \item \textbf{A2 --- Demand exceeds supply.} Dynamic pricing applies;
            the adjusted price is shown and the rider has to reconfirm before
            dispatch continues.
      \item \textbf{A3 --- Offer not accepted.} The chosen driver declines or
            does not respond in time; the system re-dispatches to the next
            candidate, or times out and informs the rider.
      \item \textbf{A4 --- Rider cancels.} The rider cancels before pickup;
            the system voids the pending trip and applies the cancellation
            policy if it applies.
    \end{itemize}
  \item[Postconditions] On success, a trip is created in the
        ``en route to pickup'' state with a driver assigned, and both parties
        are notified. On failure, no trip is created and the rider is told
        why.
\end{description}

\subsubsection*{UC-02 --- Accept and perform a trip}

\begin{description}[leftmargin=2.4cm,style=nextline,font=\normalfont\bfseries]
  \item[Primary actor] Driver.
  \item[Secondary actors] Rider; Mapping \& Geospatial Data Provider.
  \item[Preconditions] The driver is onboarded, verified, and online; the
        driver's location is being shared; a dispatch offer has been
        generated for this driver by the matching step of UC-01.
  \item[Main flow]
    \begin{enumerate}[leftmargin=1.6em,itemsep=0.15em]
      \item The system offers a nearby trip request to the driver.
      \item The driver accepts the offer.
      \item The system gives navigation to the pickup point, with rider
            status and contact.
      \item The driver reaches the pickup point and confirms pickup; the
            system starts the trip.
      \item The system gives navigation to the destination and tracks the
            trip in real time.
      \item The driver reaches the destination and ends the trip.
    \end{enumerate}
  \item[Alternative flows]
    \begin{itemize}[leftmargin=1.6em,itemsep=0.15em]
      \item \textbf{A1 --- Offer declined or ignored.} The offer expires and
            returns to UC-01 for re-dispatch to another driver.
      \item \textbf{A2 --- Rider no-show.} After a set wait, the driver
            cancels; the no-show or cancellation policy applies.
      \item \textbf{A3 --- Mid-trip change.} The rider asks for a stop or a
            new destination; the system updates route and fare accordingly.
      \item \textbf{A4 --- Trip interrupted.} The trip ends early after an
            incident; the fare is adjusted and support may step in.
    \end{itemize}
  \item[Postconditions] On success, the trip reaches the ``completed'' state
        awaiting settlement and hands off to UC-03. On decline or
        cancellation, the driver returns to the available pool.
\end{description}

\subsubsection*{UC-03 --- Complete trip and process payment}

\begin{description}[leftmargin=2.4cm,style=nextline,font=\normalfont\bfseries]
  \item[Primary actor] Rider (the billed party); \emph{triggered by} trip
        completion in UC-02.
  \item[Secondary actors] Payment Service Provider; Driver (payout
        beneficiary).
  \item[Preconditions] A trip has been completed; the rider has a registered
        payment method or an agreed cash arrangement; the fare parameters are
        known.
  \item[Main flow]
    \begin{enumerate}[leftmargin=1.6em,itemsep=0.15em]
      \item On completion, the system computes the final fare (base, distance
            and time, any dynamic pricing, tolls and fees, and promotions).
      \item The system asks the Payment Service Provider to authorize and
            capture the charge against the rider's payment method.
      \item The provider processes the charge and returns the result.
      \item The system issues a receipt to the rider and records the driver's
            earnings for the trip.
      \item The system schedules the driver payout per the payout policy.
    \end{enumerate}
  \item[Alternative flows]
    \begin{itemize}[leftmargin=1.6em,itemsep=0.15em]
      \item \textbf{A1 --- Payment declined.} The system records an
            outstanding balance, notifies the rider, and retries or asks for
            another method.
      \item \textbf{A2 --- Fare dispute.} Through UC-05, the fare is reviewed
            and, if warranted, corrected or refunded.
      \item \textbf{A3 --- Cash settlement (common in Latin American
            markets).} The fare is paid in cash directly; the system records
            the cash trip and reconciles balances instead of charging a card.
    \end{itemize}
  \item[Postconditions] On success, the trip is settled, the rider has a
        receipt, and the driver's earnings are recorded and scheduled for
        payout. On failure, the trip is flagged with an outstanding balance
        for follow-up.
\end{description}

\subsubsection*{UC-04 --- Onboard and verify a driver}

\begin{description}[leftmargin=2.4cm,style=nextline,font=\normalfont\bfseries]
  \item[Primary actor] Driver (applicant).
  \item[Secondary actors] Identity \& Background Verification Provider;
        Operations / Market Administrator.
  \item[Preconditions] The applicant has created an account; the market's
        onboarding requirements are configured.
  \item[Main flow]
    \begin{enumerate}[leftmargin=1.6em,itemsep=0.15em]
      \item The applicant submits the required identity, license, and vehicle
            information and documents.
      \item The system forwards identity verification and background
            screening to the verification provider.
      \item The provider returns the identity and screening results.
      \item The system, and where policy requires it an Operations
            Administrator, reviews the results against market requirements.
      \item On approval, the system activates the driver account, letting the
            driver go online.
    \end{enumerate}
  \item[Alternative flows]
    \begin{itemize}[leftmargin=1.6em,itemsep=0.15em]
      \item \textbf{A1 --- Disqualifying result.} The application is rejected
            or held; the applicant is notified with the reason where
            permitted.
      \item \textbf{A2 --- Incomplete documents.} The system asks for
            corrections and resubmission.
      \item \textbf{A3 --- Periodic re-screening.} An active driver is
            re-verified on a schedule; a failure can suspend the account.
    \end{itemize}
  \item[Postconditions] On approval, a verified, active driver exists and can
        accept trips (enabling UC-02). On rejection, no active driver account
        is created.
\end{description}

\subsubsection*{UC-05 --- Rate and report a trip}

\begin{description}[leftmargin=2.4cm,style=nextline,font=\normalfont\bfseries]
  \item[Primary actor] Rider or Driver (either may rate the other).
  \item[Secondary actors] Support / Safety Agent.
  \item[Preconditions] A trip has been completed (or, for a safety report, is
        in progress or just ended); the participant is authenticated.
  \item[Main flow]
    \begin{enumerate}[leftmargin=1.6em,itemsep=0.15em]
      \item After the trip, the system prompts the participant to rate the
            other party and, optionally, leave feedback.
      \item The participant submits a rating and any optional comment or tip.
      \item The system records the rating and updates the other party's
            aggregate score.
      \item If the participant files a report (safety, service, or lost
            item), the system routes it to a Support / Safety Agent.
      \item The agent reviews and resolves the report, taking the appropriate
            action per policy.
    \end{enumerate}
  \item[Alternative flows]
    \begin{itemize}[leftmargin=1.6em,itemsep=0.15em]
      \item \textbf{A1 --- No rating submitted.} The system moves on; the
            rating may be asked for again later or skipped.
      \item \textbf{A2 --- Low rating or serious report.} The trip is flagged
            for review and may lead to follow-up or account action.
      \item \textbf{A3 --- Lost item or dispute.} The case is routed to
            support with the relevant trip data for resolution.
    \end{itemize}
  \item[Postconditions] Ratings are recorded and aggregates updated; any
        report carries a tracked resolution state, and quality and safety
        signals become available for later policy.
\end{description}

%---------------------------------------------------------------------
\subsection{System Boundaries}
\label{subsec:boundaries}

This subsection sets the limits of the system: what sits inside its boundary
(the capabilities it owns and answers for), what is explicitly out of scope,
and the external interfaces through which it meets its environment. The scope
is the global Uber ride-sharing platform, in line with the worldwide figures
in Section~\ref{subsec:business-problem}; where a Latin American market
differs in a way that matters, it is noted. The test throughout is one of
control: a capability the platform designs, runs, and answers for is inside
the boundary; a capability it depends on but does not control is outside it.

\subsubsection*{Inside the boundary (in scope)}

The system owns the following:

\begin{itemize}[leftmargin=1.4em,itemsep=0.15em]
  \item Rider and driver account management, profiles, and session
        authentication.
  \item Real-time supply--demand \textbf{matching and dispatch} (candidate
        discovery, offer, and assignment).
  \item \textbf{Trip lifecycle management} across its states (requested, en
        route to pickup, in progress, completed, settled).
  \item \textbf{Dynamic pricing} (surge) computation per geographic area.
  \item \textbf{Fare calculation and settlement orchestration}: computing
        fares and coordinating charges and payouts (the payment execution
        itself is external).
  \item In-house geospatial and routing logic (ETA estimation, trip
        tracking) built on external base-map data.
  \item \textbf{Ratings, feedback, and reputation} management.
  \item Support-case management and safety-incident routing (the workflow,
        not third-party services).
  \item Notification orchestration: deciding what to send and when (message
        delivery itself is external).
  \item Marketplace configuration and operations tooling: per-market policy,
        service areas, and monitoring.
  \item Risk and fraud logic.
\end{itemize}

\noindent\emph{LATAM note.} Settlement orchestration has to handle
\textbf{cash-paid trips}, which are common in several Latin American markets;
for those trips the platform records and reconciles the fare rather than
relying on an external card charge.

\subsubsection*{Outside the boundary (out of scope)}

The following are not built or owned by the system; the platform depends on
or coexists with them:

\begin{itemize}[leftmargin=1.4em,itemsep=0.15em]
  \item Payment processing, banking, and settlement, handled by an external
        payment service provider.
  \item The actual identity verification and background screening (an
        external verification provider).
  \item Base map data and geocoding sources (an external mapping/geodata
        provider).
  \item Message-delivery infrastructure for SMS, push, and email (external
        carriers).
  \item Rider and driver devices, their operating systems, and network
        connectivity.
  \item The physical transportation itself, the vehicles, and their
        maintenance.
  \item Regulatory, tax, and compliance rule-making (the platform complies
        with rules; it does not write them).
\end{itemize}

\subsubsection*{External interfaces}

Every actor in the catalog (Table~\ref{tab:actors}) meets the platform across
the boundary, and each crossing is an integration point. The human actors
(rider, driver, support/safety agent, operations administrator) start or
service interactions with the platform. The external systems are called on by
the platform: the payment provider for authorization, capture, and payouts;
the mapping/geodata provider for base-map and geocoding data; the identity
and background verification provider for onboarding and periodic checks; and
the notification provider for message delivery. These crossings are drawn in
Figure~\ref{fig:boundary}.

\subsubsection*{System boundary diagram}

Figure~\ref{fig:boundary} shows the system as a single black box: it shows
what the platform interacts with, not how it is built inside. Opening the box
into services is left, on purpose, to the software-architecture delivery.

\begin{figure}[H]
  \centering
  \begin{tikzpicture}[
    font=\small,
    human/.style={draw, rounded corners=2pt, align=center, fill=black!4,
      minimum width=2.7cm, minimum height=0.9cm},
    system/.style={draw, rounded corners=2pt, align=center, fill=black!8,
      minimum width=2.7cm, minimum height=0.9cm},
    sysbox/.style={draw, dashed, very thick, rounded corners=6pt,
      align=center, fill=blue!5, minimum width=3.7cm, minimum height=2.6cm},
    link/.style={<->, >={Latex[length=2mm]}, semithick, black!55}
  ]
    % System-of-interest, enclosed by its (dashed) boundary
    \node[sysbox] (sys) at (0,0)
      {\textbf{Uber Ride-Sharing}\\\textbf{Platform}\\[3pt]
       \footnotesize\itshape (system-of-interest)};
    \node[anchor=south, font=\scriptsize\itshape] at (sys.north)
      {system boundary};

    % Human actors (left)
    \node[human] (rider)   at (-5, 2.1) {Rider};
    \node[human] (driver)  at (-5, 0.7) {Driver};
    \node[human] (support) at (-5,-0.7) {Support /\\Safety Agent};
    \node[human] (ops)     at (-5,-2.1) {Operations /\\Market Admin};

    % External systems (right)
    \node[system] (pay)   at (5, 2.1) {Payment\\Service Provider};
    \node[system] (map)   at (5, 0.7) {Mapping \&\\Geospatial Data};
    \node[system] (idv)   at (5,-0.7) {Identity \&\\Background Verif.};
    \node[system] (notif) at (5,-2.1) {Notification /\\Comms};

    % Interactions crossing the boundary
    \draw[link] (rider)   -- (sys);
    \draw[link] (driver)  -- (sys);
    \draw[link] (support) -- (sys);
    \draw[link] (ops)     -- (sys);
    \draw[link] (sys) -- (pay);
    \draw[link] (sys) -- (map);
    \draw[link] (sys) -- (idv);
    \draw[link] (sys) -- (notif);

    % Column captions
    \node[font=\scriptsize\bfseries] at (-5, 3.0) {Human actors};
    \node[font=\scriptsize\bfseries] at ( 5, 3.0) {External systems};
  \end{tikzpicture}
  \caption{System boundary (context) diagram for the global Uber
    ride-sharing platform. The platform is shown as a black box; arrows are
    interactions crossing the boundary, not internal structure.}
  \label{fig:boundary}
\end{figure}